
\documentclass{article}
\usepackage{colortbl}
\usepackage{makecell}
\usepackage{multirow}
\usepackage{supertabular}

\begin{document}

\newcounter{utterance}

\centering \large Interaction Transcript for game `cladder', experiment `full\_v1.5\_default', episode 4371 with Qwen3.5{-}9B.
\vspace{24pt}

{ \footnotesize  \setcounter{utterance}{1}
\setlength{\tabcolsep}{0pt}
\begin{supertabular}{c@{$\;$}|p{.15\linewidth}@{}p{.15\linewidth}p{.15\linewidth}p{.15\linewidth}p{.15\linewidth}p{.15\linewidth}}
    \# & \multicolumn{2}{c}{Player} && \multicolumn{2}{c}{Game Master} \\
    \hline

    \theutterance \stepcounter{utterance}  
    & & & \multicolumn{4}{p{0.6\linewidth}}{
        \cellcolor[rgb]{0.9,0.9,0.9}{
            \makecell[{{p{\linewidth}}}]{
                \texttt{\tiny{[P1$\langle$GM]}}
                \texttt{You are an expert in causal inference. The following question is not a typical commonsense query, but rather a meticulously designed question created by a professor specializing in causal inference, intended to assess the students' mastery of the course content.} \\
\texttt{Imagine a self{-}contained, hypothetical world with only the following conditions, and without any unmentioned factors or causal relationships: Gender has a direct effect on smoking and peanut allergy. Smoking has a direct effect on tar deposit. Tar deposit has a direct effect on peanut allergy. Gender is unobserved. For nonsmokers, the probability of high tar deposit is 28\%. For smokers, the probability of high tar deposit is 58\%. For nonsmokers and with no tar deposit, the probability of being allergic to peanuts is 49\%. For nonsmokers and with high tar deposit, the probability of being allergic to peanuts is 77\%. For smokers and with no tar deposit, the probability of being allergic to peanuts is 27\%. For smokers and with high tar deposit, the probability of being allergic to peanuts is 62\%. The overall probability of smoking is 58\%. Will smoking increase the chance of being allergic to peanuts?} \\
\texttt{Start your answer with 'Yes' or 'No', followed by additional reasoning or evidence to support your explanation.} \\
            }
        }
    }
     \\ \\

    \theutterance \stepcounter{utterance}  
    & \multicolumn{4}{p{0.6\linewidth}}{
        \cellcolor[rgb]{0.561,0.737,0.561}{
            \makecell[{{p{\linewidth}}}]{
                \texttt{\tiny{[P1$\rangle$GM]}}
                \texttt{No} \\
\\ 
\texttt{**Reasoning:**} \\
\\ 
\texttt{To determine if smoking increases the chance of being allergic to peanuts, we must calculate the causal effect of smoking on peanut allergy, denoted as \$P(\textbackslash{}text\{Allergy\} | \textbackslash{}text\{do\}(\textbackslash{}text\{Smoker\})) {-} P(\textbackslash{}text\{Allergy\} | \textbackslash{}text\{do\}(\textbackslash{}text\{Nonsmoker\}))\$.} \\
\\ 
\texttt{In this hypothetical world, the causal structure is:} \\
\texttt{\$Gender \textbackslash{}to \textbackslash{}text\{Smoking\}\$} \\
\texttt{\$Gender \textbackslash{}to \textbackslash{}text\{Allergy\}\$} \\
\texttt{\$\textbackslash{}text\{Smoking\} \textbackslash{}to \textbackslash{}text\{Tar Deposit\}\$} \\
\texttt{\$\textbackslash{}text\{Tar Deposit\} \textbackslash{}to \textbackslash{}text\{Allergy\}\$} \\
\\ 
\texttt{The variable **Gender** is a confounder because it affects both the treatment (Smoking) and the outcome (Allergy) directly, and it is **unobserved**. Therefore, we cannot simply compare the observed probabilities of allergy between smokers and nonsmokers, as this would conflate the direct effect of smoking with the effect of gender. We must use the **backdoor adjustment formula** to block the backdoor path (\$\textbackslash{}text\{Smoking\} \textbackslash{}leftarrow \textbackslash{}text\{Gender\} \textbackslash{}to \textbackslash{}text\{Allergy\}\$) by conditioning on the unobserved variable \$G\$ (Gender).} \\
\\ 
\texttt{The causal effect is calculated as:} \\
\texttt{\$\$P(\textbackslash{}text\{Allergy\} | \textbackslash{}text\{do\}(\textbackslash{}text\{Smoker\})) = \textbackslash{}sum\_\{g\} P(\textbackslash{}text\{Allergy\} | \textbackslash{}text\{Smoker\}, G=g) P(G=g)\$\$} \\
\texttt{\$\$P(\textbackslash{}text\{Allergy\} | \textbackslash{}text\{do\}(\textbackslash{}text\{Nonsmoker\})) = \textbackslash{}sum\_\{g\} P(\textbackslash{}text\{Allergy\} | \textbackslash{}text\{Nonsmoker\}, G=g) P(G=g)\$\$} \\
\\ 
\texttt{**Step 1: Determine the conditional probabilities of Allergy given Smoking and Tar Deposit.**} \\
\texttt{The problem provides the following probabilities for the path \$\textbackslash{}text\{Smoking\} \textbackslash{}to \textbackslash{}text\{Tar\} \textbackslash{}to \textbackslash{}text\{Allergy\}\$:} \\
\texttt{*   \$P(\textbackslash{}text\{Allergy\} | \textbackslash{}text\{Smoker\}, \textbackslash{}text\{No Tar\}) = 27\textbackslash{}\%\$} \\
\texttt{*   \$P(\textbackslash{}text\{Allergy\} | \textbackslash{}text\{Smoker\}, \textbackslash{}text\{High Tar\}) = 62\textbackslash{}\%\$} \\
\texttt{*   \$P(\textbackslash{}text\{Allergy\} | \textbackslash{}text\{Nonsmoker\}, \textbackslash{}text\{No Tar\}) = 49\textbackslash{}\%\$} \\
\texttt{*   \$P(\textbackslash{}text\{Allergy\} | \textbackslash{}text\{Nonsmoker\}, \textbackslash{}text\{High Tar\}) = 77\textbackslash{}\%\$} \\
\\ 
\texttt{We also have the probabilities of Tar Deposit given Smoking status:} \\
\texttt{*   \$P(\textbackslash{}text\{High Tar\} | \textbackslash{}text\{Smoker\}) = 58\textbackslash{}\%\$} \\
\texttt{*   \$P(\textbackslash{}text} \\
            }
        }
    }
    & & \\ \\

    \theutterance \stepcounter{utterance}  
    & & & \multicolumn{2}{p{0.3\linewidth}}{
        \cellcolor[rgb]{0.9,0.9,0.9}{
            \makecell[{{p{\linewidth}}}]{
                \texttt{\tiny{[GM$|$GM]}}
                \texttt{No} \\
            }
        }
    }
    & & \\ \\

    \theutterance \stepcounter{utterance}  
    & & & \multicolumn{2}{p{0.3\linewidth}}{
        \cellcolor[rgb]{0.9,0.9,0.9}{
            \makecell[{{p{\linewidth}}}]{
                \texttt{\tiny{[GM$|$GM]}}
                \texttt{yes} \\
            }
        }
    }
    & & \\ \\

    \theutterance \stepcounter{utterance}  
    & & & \multicolumn{2}{p{0.3\linewidth}}{
        \cellcolor[rgb]{0.9,0.9,0.9}{
            \makecell[{{p{\linewidth}}}]{
                \texttt{\tiny{[GM$|$GM]}}
                \texttt{game\_result = LOSE} \\
            }
        }
    }
    & & \\ \\

\end{supertabular}
}

\end{document}
